\section{Lecture 5: Generalization and Overfitting}

\subsection{Motivation}

In previous lectures, we formulated learning as:
\begin{itemize}
    \item approximation of an unknown function,
    \item estimation from random data,
    \item optimization of an objective function.
\end{itemize}

\medskip

\noindent
However, a central challenge remains:

\medskip

\noindent
\textit{Why should a model that fits the training data perform well on unseen data?}

\medskip

\noindent
This question lies at the heart of machine learning and is addressed by the concept of \emph{generalization}.

\medskip

\noindent
\textbf{Guiding idea.}  
Learning is successful only if performance extends beyond the training data.

\subsection{Training and Test Performance}

Let $\mathcal{D}_{\text{train}}$ denote the training dataset and $\mathcal{D}_{\text{test}}$ denote unseen data drawn from the same distribution.

\medskip

\noindent
\textbf{Training error.}
\[
\hat{R}_{\text{train}}(f) = \frac{1}{n} \sum_{(x_i,y_i)\in \mathcal{D}_{\text{train}}} \ell(f(x_i), y_i).
\]

\medskip

\noindent
\textbf{Test error.}
\[
\hat{R}_{\text{test}}(f) = \frac{1}{m} \sum_{(x_j,y_j)\in \mathcal{D}_{\text{test}}} \ell(f(x_j), y_j).
\]

\medskip

\noindent
\textbf{True risk.}
\[
R(f) = \mathbb{E}_{(x,y)\sim P}[\ell(f(x), y)].
\]

\medskip

\noindent
\textbf{Generalization gap.}
\[
\hat{R}_{\text{test}}(f) - \hat{R}_{\text{train}}(f).
\]

\medskip

\noindent
\textbf{Interpretation.}
\begin{itemize}
    \item Training error measures fit to observed data
    \item Test error measures predictive performance
    \item The gap reflects overfitting
\end{itemize}

\medskip

\noindent
\textbf{Insight.}  
Low training error alone is not sufficient; generalization is the true objective.

\subsection{Overfitting and Underfitting}

\textbf{Overfitting.}
\begin{itemize}
    \item Model fits training data extremely well
    \item Captures noise rather than underlying structure
    \item High variance, poor test performance
\end{itemize}

\medskip

\noindent
\textbf{Underfitting.}
\begin{itemize}
    \item Model is too simple
    \item Fails to capture underlying patterns
    \item High bias, poor performance on both training and test data
\end{itemize}

\medskip

\noindent
\textbf{Typical behavior.}
\begin{itemize}
    \item Increasing model complexity decreases training error
    \item Test error decreases initially, then increases (classical view)
\end{itemize}

\medskip

\noindent
\textbf{Illustration.}
\begin{itemize}
    \item Simple model: smooth but inaccurate
    \item Complex model: oscillatory, fits noise
\end{itemize}

\medskip

\noindent
\textbf{Insight.}  
There exists an optimal level of complexity balancing bias and variance.

\subsection{Model Capacity and Complexity}

\textbf{Definition.}  
Model capacity refers to the ability of a hypothesis space to fit a wide range of functions.

\medskip

\noindent
\textbf{Examples.}
\begin{itemize}
    \item Linear models: low capacity
    \item High-degree polynomials: high capacity
    \item Deep neural networks: extremely high capacity
\end{itemize}

\medskip

\noindent
\textbf{Measures of capacity.}
\begin{itemize}
    \item Number of parameters
    \item Flexibility of function class
    \item Ability to fit arbitrary labels
\end{itemize}

\medskip

\noindent
\textbf{Interpolation phenomenon.}  
High-capacity models can fit training data exactly:
\[
\hat{R}_{\text{train}}(f) = 0.
\]

\medskip

\noindent
\textbf{Classical view.}
\begin{itemize}
    \item High capacity $\Rightarrow$ overfitting
\end{itemize}

\medskip

\noindent
\textbf{Modern observation.}
\begin{itemize}
    \item Overparameterized models can still generalize well
\end{itemize}

\medskip

\noindent
\textbf{Insight.}  
Capacity alone does not determine generalization; the interaction with optimization is crucial.

\subsection{Informal Generalization Principles}

Although a complete theory of generalization is complex, several guiding principles are useful.

\medskip

\noindent
\textbf{1. Simplicity (Occam's principle).}
\begin{itemize}
    \item Simpler models tend to generalize better
    \item Regularization promotes simplicity
\end{itemize}

\medskip

\noindent
\textbf{2. Stability.}
\begin{itemize}
    \item Models that are insensitive to small data changes generalize better
\end{itemize}

\medskip

\noindent
\textbf{3. Margin and robustness.}
\begin{itemize}
    \item Larger margins in classification often improve generalization
\end{itemize}

\medskip

\noindent
\textbf{4. Implicit bias of optimization.}
\begin{itemize}
    \item Optimization algorithms favor certain solutions
    \item Gradient descent often finds structured or “simple” solutions
\end{itemize}

\medskip

\noindent
\textbf{5. Data quality and representation.}
\begin{itemize}
    \item Better features lead to better generalization
\end{itemize}

\medskip

\noindent
\textbf{Insight.}  
Generalization depends on a combination of model structure, data, and optimization dynamics.

\subsection{A Unifying Perspective}

Generalization can be understood as the relationship between:

\begin{itemize}
    \item \textbf{Data:} finite samples from a distribution
    \item \textbf{Model:} hypothesis space and capacity
    \item \textbf{Optimization:} how solutions are selected
\end{itemize}

\medskip

\noindent
\textbf{Core idea.}  
Learning selects one solution among many that fit the data; generalization depends on which solution is chosen.

\medskip

\noindent
\textbf{Key insight.}  
The generalization behavior of modern machine learning systems cannot be explained by model capacity alone.

\subsection{Outlook}

In this lecture, we studied:

\begin{itemize}
    \item training and test performance,
    \item overfitting and underfitting,
    \item model capacity and complexity,
    \item informal principles of generalization.
\end{itemize}

\medskip

\noindent
These ideas motivate the need for methods that control complexity and improve generalization.

\medskip

\noindent
In the next lecture, we study \emph{regularization}, which introduces explicit constraints and penalties to guide learning.

\medskip

\noindent
\textbf{Guiding principle.}  
The goal of learning is not to fit the data perfectly, but to capture the underlying structure that generates it.